% !TeX program = xelatex
\documentclass[11pt,a4paper]{article}
\usepackage{fontspec}
\IfFontExistsTF{Times New Roman}{\setmainfont{Times New Roman}}{\setmainfont{TeX Gyre Termes}}
\IfFontExistsTF{Consolas}{\setmonofont{Consolas}[Scale=0.84]}{\setmonofont{Latin Modern Mono}}
\usepackage[margin=2cm,headheight=15pt]{geometry}
\usepackage{booktabs,tabularx,array,fancyhdr,enumitem,amssymb}
\usepackage[table]{xcolor}
\usepackage[unicode,colorlinks=true,urlcolor=blue]{hyperref}
\hypersetup{pdftitle={Hệ thống xử lý video đa camera},pdfauthor={Trần Danh Thiện}}
\definecolor{navy}{HTML}{173B61}
\definecolor{lightblue}{HTML}{E6F0F8}
\setlength{\parindent}{0pt}
\setlength{\parskip}{5pt}
\renewcommand{\arraystretch}{1.17}
\setlist[itemize]{leftmargin=16pt,itemsep=2pt,topsep=3pt}
\pagestyle{fancy}\fancyhf{}
\fancyhead[L]{\small Trần Danh Thiện | Xử lý video đa camera}
\fancyhead[R]{\small 26/10/2026}
\fancyfoot[C]{\thepage}
\newcolumntype{Y}{>{\raggedright\arraybackslash}X}
\newcommand{\code}[1]{\texttt{\detokenize{#1}}}
\newcommand{\note}[1]{\par\smallskip\colorbox{lightblue}{\parbox{\dimexpr\linewidth-2\fboxsep\relax}{#1}}\par\smallskip}
\begin{document}
\begin{center}
{\fontsize{21}{25}\selectfont\color{navy}\bfseries BÁO CÁO KỸ THUẬT}\\[6pt]
{\fontsize{16}{20}\selectfont\color{navy}\bfseries HỆ THỐNG XỬ LÝ VIDEO ĐA CAMERA}\\[8pt]
\textbf{Người thực hiện: Trần Danh Thiện}\\
\href{https://github.com/thien2603/Detect-and-Tracking}{github.com/thien2603/Detect-and-Tracking}\\
{\small Cập nhật 26/10/2026; benchmark đo ngày 23/10/2026.}
\end{center}
\note{\textbf{Kết quả chính:} chạy đồng thời 4 nguồn, nhận diện người/xe, tracking, đếm hai chiều và cảnh báo vùng. Microbatch đạt trung bình \textbf{11,07 FPS ảnh mới/camera}; độ trễ lớn nhất \textbf{1129 ms}, chỉ 3/4 camera đạt dưới 1 giây. Chưa đạt toàn bộ tiêu chí hiệu năng.}

\section{Kiến trúc và luồng hoạt động}
Bốn nguồn gồm ba góc Bellevue và một góc Crosswalk có người đi bộ. Khi benchmark, file video được phát lại qua RTSP; mỗi camera có tracker, bộ đếm và đường RTSP đầu ra riêng.
\begin{center}
\fcolorbox{navy}{lightblue}{\parbox{0.92\linewidth}{\centering
Nguồn $\rightarrow$ Capture giữ ảnh mới nhất $\rightarrow$ Scheduler $\rightarrow$ YOLO11n\\[3pt]
$\downarrow$\\
Tracker từng camera $\rightarrow$ Quy tắc vạch / polygon\\
$\swarrow$\hspace{40mm}$\searrow$\\
SQLite + snapshot\hspace{12mm}Vẽ kết quả $\rightarrow$ FFmpeg $\rightarrow$ RTSP}}
\end{center}
Capture độc lập giúp một nguồn lỗi không kéo các nguồn khác dừng theo. Scheduler giữ tối đa một yêu cầu chờ/camera; ảnh cũ được thay bằng ảnh mới. Ghi bằng chứng và phát RTSP tách khỏi suy luận. FastAPI cung cấp trạng thái, FPS, bật/tắt camera và truy vấn sự kiện.

{\small\begin{tabularx}{\linewidth}{p{64mm}Y}
\toprule\rowcolor{lightblue}\textbf{Nhóm mã nguồn} & \textbf{Chức năng}\\\midrule
\path{app/camera/capture.py} & Đọc file/RTSP/webcam, latest-frame, reconnect.\\
\path{app/inference/} & YOLO, warmup và điều phối batch.\\
\path{app/inference/tracker.py}; \path{motion.py} & IoU, Kalman và ID riêng theo camera.\\
\path{app/pipeline.py}; \path{app/events/} & Quy tắc đếm/vùng; SQLite, snapshot.\\
\path{app/streaming/}; \path{app/manager.py} & Overlay, FFmpeg và vòng đời camera.\\
\path{app/main.py}; \path{configs/} & API và cấu hình YAML.\\
\bottomrule\end{tabularx}}

\section{Thuật toán và quy tắc sự kiện}
\textbf{Detection/tracking:} YOLO11n tìm bbox và lớp đối tượng. Tracker ghép bbox cùng lớp bằng IoU; bản tối ưu dùng Kalman dự đoán vị trí trước khi ghép, hỗ trợ chuyển động nhanh. Che khuất hoặc thiếu detection lâu vẫn có thể gây đổi ID.

\textbf{Đếm qua vạch:} lấy tâm cạnh dưới bbox; tăng bộ đếm theo chiều khi điểm đổi phía và cắt đoạn vạch hữu hạn, có deadband/cooldown.

\textbf{Vào vùng:} kiểm tra điểm chân hộp trong polygon; ghi sự kiện khi đi vào, không ghi lại mỗi frame nếu đứng trong vùng. Lý do \code{first_seen_inside} được lưu riêng khi lần đầu thấy đã ở trong vùng. Retry cùng event ID không chèn trùng vào SQLite.

\note{\textbf{Giới hạn:} dự đoán chỉ để vẽ; sự kiện dùng bbox YOLO đo thật. Hộp quá tuổi 700 ms bị ẩn. Chưa có ground truth để công bố accuracy; đổi track ID có thể gây đếm lại cùng vật thể.}

\clearpage
\section{Cấu hình và yêu cầu chức năng}
Cấu hình chính: \path{configs/submission-optimized.yaml}. Vạch/polygon được căn thủ công theo từng góc camera cố định, lưu tọa độ pixel trong YAML. Vùng xanh khoảng $3/4$ và đỏ khoảng $1/4$ ROI là vùng minh họa, chưa quy đổi thành diện tích thực.

{\small\begin{tabularx}{\linewidth}{p{42mm}Y}
\toprule\rowcolor{lightblue}\textbf{Tham số} & \textbf{Giá trị trong ma trận 1/4/8 camera}\\\midrule
Model / thiết bị & YOLO11n, CUDA RTX 3050 Laptop, FP16.\\
Ảnh / confidence & \code{imgsz=480}, ảnh vuông; confidence 0,25.\\
Lớp & person, car, motorcycle, bus, truck.\\
Điều phối & Serial: batch 1; microbatch: tối đa 4, chờ 15 ms.\\
Render / detect & Bất đồng bộ; \code{detect_interval=2}; Kalman/IoU.\\
RTSP & H.264 NVENC, mục tiêu 14 FPS, tối đa $960\times540$, giữ tỷ lệ.\\
Vòng đời & Hàng đợi hữu hạn, bỏ ảnh quá cũ, tự khởi động lại FFmpeg.\\
\bottomrule\end{tabularx}}

{\small\begin{tabularx}{\linewidth}{p{9mm}p{63mm}Y}
\toprule\rowcolor{lightblue}\textbf{Mã} & \textbf{Yêu cầu đề bài} & \textbf{Kết quả đối chiếu}\\\midrule
F01 & File/RTSP/webcam; reconnect; cô lập lỗi & File/RTSP đã chạy; reconnect có test tích hợp. Webcam có code, chưa thử thiết bị thật.\\
F02 & Luôn xử lý frame mới nhất & Có latest-frame slot, queue hữu hạn và test thay yêu cầu cũ.\\
F03 & YOLO v8+, người/xe; class/conf/imgsz từ config & YOLO11n; có người/xe thật trong nguồn; tham số YAML.\\
F04 & Tracking & ID theo camera; test cùng lớp, mất quan sát ngắn và chuyển động nhanh.\\
F05 & Đếm qua vạch hai chiều & Unit test vạch hữu hạn, deadband/cooldown; demo có line crossing.\\
F06 & Phát hiện vào polygon & Test biên, entry/stay/exit/reentry; demo có intrusion.\\
F07 & Một event/lần vi phạm, SQLite và ảnh & Có UTC, camera, track, loại và snapshot; retry không chèn trùng.\\
F08 & Overlay, RTSP riêng, VLC, FFmpeg restart & RTSP decode thật, overlay và restart có kiểm thử; chưa nghiệm thu VLC GUI.\\
F09 & YAML; fail-fast; thêm camera bằng config & Schema/preflight trước khi chạy; test cấu hình sai và thêm nguồn.\\
F10 & FastAPI trạng thái/FPS, bật/tắt, lọc event & Có endpoint và test API, lọc camera/thời gian, snapshot.\\
\bottomrule\end{tabularx}}

\textbf{API chính:} \code{GET /api/v1/cameras}; \code{PUT /api/v1/cameras/{id}/enabled}; \code{GET /api/v1/events}. Tài liệu tương tác tại \code{/docs} khi service đang chạy.

\textbf{Điểm cần phân biệt:} có đầy đủ các nhóm chức năng trong mã nguồn, nhưng F01 còn thiếu nghiệm thu webcam vật lý, F08 còn thiếu mở bằng VLC. Độ chính xác đếm/tracking chưa được đo bằng nhãn chuẩn; đề không đặt ngưỡng accuracy cụ thể.

\clearpage
\section{Benchmark và lựa chọn kiến trúc}
\textbf{Máy đo:} Lenovo Legion 5 15ARH7; Ryzen 5 6600H (6 nhân/12 luồng); RAM lắp 8 GB, dùng được khoảng 7,19 GiB; RTX 3050 Laptop 4 GB, driver 546.92. Windows build 26200, Python 3.10.11, PyTorch CUDA, MediaMTX 1.9.3.

Mỗi phương án đo 60 giây sau 8 giây warmup. Tám camera là tám luồng logic lấy từ bốn video. \textbf{Detect FPS} đếm ảnh YOLO hoàn tất; \textbf{unique FPS} đếm timestamp ảnh nguồn khác nhau giải mã từ RTSP đầu ra. Render bất đồng bộ giúp ảnh xem mới hơn nhịp detect.

Độ trễ đo từ replay tới client giải mã RTSP trên cùng máy; chưa gồm hiển thị VLC và phơi sáng camera vật lý. Cột trễ đạt dùng max $<1000$ ms và marker coverage $\geq95\%$. FPS là trung bình camera; RAM là peak RSS cộng các tiến trình, có thể tính trùng trang dùng chung.

{\small\begin{tabularx}{\linewidth}{r Y r r r c r}
\toprule\rowcolor{lightblue}\textbf{Cam} & \textbf{Cách chạy} & \textbf{Detect} & \textbf{Unique} & \textbf{Max ms} & \textbf{Trễ đạt} & \textbf{RAM MiB}\\\midrule
1 & Serial & 6,48 & 12,55 & 639 & 1/1 & 1555\\
1 & Microbatch & 6,52 & 12,38 & 723 & 1/1 & 2036\\
4 & Serial & 1,14 & 10,52 & 1091 & 1/4 & 1964\\
4 & Microbatch & 3,05 & 11,07 & 1129 & 3/4 & 2583\\
8 & Serial & 0,17 & 0,91 & 10191 & 0/8 & 1426\\
8 & Microbatch & 0,37 & 0,98 & 7277 & 0/8 & 1582\\
\bottomrule\end{tabularx}}

\textbf{Nhận xét:} 1 camera đạt FPS và độ trễ. 4 camera đạt unique FPS $\geq10$ cho từng camera, nhưng chưa đạt độ trễ cho tất cả. 8 camera quá tải; một số luồng gặp lỗi encoder NVENC và restart, nên kết quả phản ánh cả giới hạn xử lý/phát video.

\textbf{Chọn shared microbatch:} ở 4 camera, detect tăng từ 1,14 lên 3,05 FPS/camera; CPU tổng giảm từ 702\% xuống 646\%; số camera đạt độ trễ tăng từ 1/4 lên 3/4. CPU cộng theo tiến trình có thể vượt 100\% trên máy nhiều nhân. Max 1129 ms vẫn chưa đạt ngưỡng đề.

Đối chứng riêng cùng 4 nguồn, FP32 và libx264 cho bảng dưới. Cấu hình này khác bảng trên, chỉ dùng so sánh số model trong cùng phép thử.

{\small\begin{tabularx}{\linewidth}{Y r r r}
\toprule\rowcolor{lightblue}\textbf{Phương án} & \textbf{Detect FPS/cam} & \textbf{CPU \%} & \textbf{RSS MiB}\\\midrule
Một model chung / batch & 5,01 & 576 & 1937\\
Bốn model riêng & 4,82 & 675 & 1991\\
Bốn model + CUDA stream riêng & 4,94 & 676 & 2008\\
\bottomrule\end{tabularx}}
Bốn model trên cùng GPU chưa cải thiện tốc độ; vì vậy giữ model chung. Mỗi cấu hình mới đo một lượt, chưa có kiểm tra dài hạn hoặc khoảng tin cậy.

\textbf{Bằng chứng:} hướng dẫn và số liệu trong repository:\\
{\small\path{docs/benchmark-rerun-2026-10-23.md}\\
\path{docs/model-modes-2026-10-23.md}}\par
JSON thô giữ local (thư mục data không đưa lên GitHub):\\
{\small\path{data/benchmark-rerun-20261023-01/}\\
\path{data/benchmark-rerun-20261023-4-8-02/}}

\clearpage
\section{Kiểm thử và đối chiếu nghiệm thu}
Test kiểm tra tracking, vạch/vùng, event, API và vòng đời tài nguyên. Integration dùng FFmpeg/MediaMTX thật, ngắt/nối lại nguồn và dừng FFmpeg để thử restart. Detector giả dùng trong test hạ tầng; benchmark hiệu năng dùng YOLO thật.

\textbf{Kiểm tra ngày 26/10/2026:} \textbf{94 test chính và 7 subtest đều pass}, không lỗi/bỏ qua, trong 89,10 giây. JUnit ghi tại \path{tmp/criteria-pytest-complete-20261026.xml}. Ruff đã qua; Docker Compose parse hợp lệ. Cần môi trường có dependency dev và đường dẫn FFmpeg/MediaMTX để chạy đủ test.

{\small\begin{tabularx}{\linewidth}{p{58mm}p{27mm}Y}
\toprule\rowcolor{lightblue}\textbf{Tiêu chí} & \textbf{Kết luận} & \textbf{Căn cứ / phần còn thiếu}\\\midrule
4 camera, mỗi camera $\geq10$ FPS & Đạt đầu ra & Benchmark 4/4 theo unique FPS; demo 120 giây.\\
Trễ đầu-cuối $<1$ giây & Chưa đạt đủ & Microbatch 3/4; max 1129 ms. Chưa đo tới VLC.\\
Ghi rõ máy và cách đo & Đạt & Phần cứng, cấu hình, phạm vi đo ở trang 3.\\
Ctrl+C dừng sạch & Có bằng chứng & Cleanup các lượt hoàn tất: thoát mềm, không còn tiến trình con.\\
Logging file, FPS/trễ/reconnect & Đạt mã & Logging xoay vòng, số liệu theo camera.\\
Unit test vạch và polygon & Có kiểm thử & Test biên, hai chiều, cooldown, entry/reentry.\\
Sử dụng Docker & Chưa nghiệm thu & Có Dockerfile/Compose, Engine đã chạy; chưa build/up stack vì thiếu RAM.\\
Đo 1/4/8, ít nhất hai cách & Đạt đo đạc & Sáu tổ hợp serial/microbatch, 60 giây/lượt. Đề không yêu cầu 8 camera đạt 10 FPS.\\
Webcam và mở RTSP bằng VLC & Còn xác minh & Có code webcam và RTSP decode thật; thiếu thiết bị/VLC.\\
\bottomrule\end{tabularx}}

\section{Hồ sơ nộp và việc còn lại}
\textbf{Mã nguồn:} đã fetch GitHub, xác nhận \code{main} local trùng \code{origin/main} tại commit \code{743923d} trước lần cập nhật tài liệu này. Repository lưu source/config/test; model và video tải riêng.

\textbf{Demo:} \code{Demo_TranDanhThien.mp4}, khoảng 120 giây, ghép bốn nguồn đang chạy đồng thời. Có 391 sự kiện trong khoảng ghi hình: 319 intrusion, 56 region entry, 16 line crossing. Số sự kiện không tương đương số đối tượng duy nhất. Đã giải mã toàn bộ clip để kiểm tra.

\textbf{Việc chưa hoàn tất:}
\begin{itemize}
\item Giảm độ trễ camera còn lại xuống dưới 1000 ms và đo lại cùng điều kiện.
\item Build/up Docker project; kiểm tra API, bốn RTSP, event và shutdown container.
\item Thử webcam thật; mở bốn RTSP đầu ra bằng VLC.
\item Upload PDF/MP4 vào thư mục Drive \code{TranDanhThien}, kiểm tra quyền xem.
\end{itemize}
\note{\textbf{Kết luận:} chức năng chính và benchmark đã có bằng chứng. Chưa đủ căn cứ xác nhận đạt toàn bộ đề vì còn độ trễ 4 camera, Docker project và các bước nghiệm thu/nộp bài nêu trên.}
\end{document}
